\chapter{Discussion}
\label{ch:discussion}

Les chapitres précédents ont décrit deux manières de résoudre la même tâche, puis les ont mesurées sur le même banc (chapitre~\ref{ch:evaluation}). Ce chapitre prend du recul. Plutôt que de dresser un catalogue de fonctionnalités, nous revenons sur ce que la mise en {\oe}uvre a réellement \emph{révélé} : les limites que chaque module a fini par exposer, les choix techniques qui se sont imposés une fois confrontés au matériel, et ce que la contrainte de \emph{sous-actionnement} du bras a dicté à toute la stratégie de saisie. Plusieurs de ces points ne sont pas des décisions prises \emph{a priori} : ce sont des conclusions tirées d'essais qui ont d'abord échoué. Nous nous efforçons de les restituer dans cet ordre --- hypothèse, essai, problème observé, correctif ---, car c'est là que s'est concentré l'essentiel de l'apprentissage du projet. Nous terminons par les limites que nous avons \emph{choisies} --- l'évitement d'obstacles et l'occlusion --- puis par les pistes qui prolongeraient le travail.

\section{Synthèse des limites identifiées}
\label{sec:disc-limites}

Le pipeline complet étant posé, ses points faibles se laissent regrouper. Chacun a été rencontré \emph{en pratique}, puis corrigé autant que possible ; les chapitres techniques les ont exposés en détail, et nous les rassemblons ici pour en dégager le trait commun.

\paragraph{La détection : deux angles morts opposés.}
Les deux détecteurs se paient d'une faiblesse complémentaire (chapitre~\ref{ch:perception}). Le seuillage HSV ne connaît que la couleur : il confond le cube orange avec les parties oranges du robot, s'égare sur les teintes \emph{achromatiques} (le blanc et le noir n'ont pas de teinte définie), et se recale à chaque changement d'éclairage. Le détecteur appris OWL-ViTv2 lève ces confusions par le \emph{sens} --- forme et contexte plutôt que teinte ---, mais au prix d'une inférence \emph{lente} (quelques secondes par image, d'où une détection ponctuelle et non un suivi continu) et d'une simple \emph{boîte englobante} sans masque : un contour trop grossier pour fixer finement l'orientation d'un objet rond, tel le cylindre. Aucun des deux ne réunit rapidité, robustesse sémantique et contour précis --- un manque que les pistes reprendront (section~\ref{sec:disc-pistes}).

\paragraph{La précision géométrique : la hauteur est l'axe faible.}
La calibration \emph{{\oe}il-main} laisse des résidus de l'ordre du demi-centimètre (chapitre~\ref{ch:perception}, table~\ref{tab:handeye-valid}), qui se \emph{propagent} : l'objet est localisé, puis la pince posée, à quelques millimètres près. Le point à retenir n'est pas ce chiffre global, mais son \emph{anisotropie} --- l'axe le moins fiable est la \emph{hauteur}. La paire fixe surplombe la table, si bien que sa profondeur bruitée, direction la plus incertaine de toute triangulation stéréo (chapitre~\ref{ch:perception}), se projette pour l'essentiel sur la verticale. Or c'est justement l'axe que le raffinement \texttt{cam\_2} \emph{ne corrige pas} : il ne rattrape que le plan $(x, y)$, la hauteur restant celle de la stéréo (chapitre~\ref{ch:controle}). Deux marges en limitent les effets --- l'ouverture laissée large de $10$ mm de chaque côté (chapitre~\ref{ch:planification}), et la correction \texttt{cam\_2} qui efface en moyenne $13$ mm d'erreur latérale avant la descente (chapitre~\ref{ch:evaluation}) ---, mais pour un objet fin, cette faiblesse en hauteur reste le vrai plancher de fiabilité.

\paragraph{Un point non résolu : le roulis en prise inclinée.}
En prise \emph{inclinée} ($45^{\circ}$), l'orientation du roulis de poignet (\texttt{wrist\_roll}) est moins fiable qu'à la verticale. Une part en est comprise --- la convention de roulis retenue ignore le roulis \emph{réel} de l'objet, d'où une légère asymétrie gauche/droite (chapitre~\ref{ch:planification}). Mais nous avons aussi observé, à certaines positions, un comportement du roulis plus \emph{erratique} que cette seule asymétrie n'explique, dont nous n'avons pas établi la cause avec certitude. Nous le consignons comme une limite \emph{ouverte} du contrôle en régime incliné plutôt que d'en risquer une explication non vérifiée --- ce régime incliné étant, rappelons-le, celui où se concentrent déjà les échecs de loin (chapitre~\ref{ch:evaluation}).

\section{Choix discutables et sous-actionnement}
\label{sec:disc-choix}

Toute limite de perception peut, en principe, se corriger par un meilleur modèle. D'autres choix, en revanche, tiennent à la \emph{cinématique} du bras ou à des arbitrages assumés en cours de projet. Cette section revient sur les plus discutables --- à commencer par la manière dont l'\emph{angle de saisie} est décidé, qui découle directement du sous-actionnement introduit au chapitre~\ref{ch:planification}.

\paragraph{L'angle d'attaque : la position, non la forme.}
C'est, rétrospectivement, le choix le plus discutable de la planification de saisie. Le chapitre~\ref{ch:planification} l'a posé : l'\emph{ouverture} et l'\emph{orientation} des mâchoires s'adaptent à la géométrie mesurée, mais l'\emph{angle d'attaque}, lui, est dicté par la \emph{position} de l'objet --- sa distance, la hauteur de son sommet --- et son atteignabilité, non par sa forme. Nous sommes d'abord partis d'un cadre purement \emph{top-down} --- la pince à la verticale ---, puis l'avons généralisé en un \emph{balayage} d'inclinaisons ($0^{\circ}$, $45^{\circ}$, $90^{\circ}$) sélectionnées par la position ; ce n'est qu'ensuite que nous avons pleinement mesuré qu'une saisie \emph{réellement} adaptative choisirait l'angle d'après la \emph{forme et la pose} de l'objet. À position et encombrement égaux, en effet, un cube, un cylindre ou une balle reçoivent aujourd'hui le \emph{même} angle d'attaque, quelle que soit leur forme.

Plus fondamentalement, le planificateur raisonne sur une \emph{boîte englobante} --- position, encombrement, axe principal ---, non sur la \emph{forme réelle} de l'objet ni sur ses \emph{prises naturelles}. Une tasse creuse le montre bien --- un objet hors de nos primitives, que notre chaîne n'a pas eu à saisir. Le planificateur la traiterait comme une boîte : il ouvrirait les mâchoires à sa largeur et les refermerait sur ses \emph{parois}. Cela \emph{fonctionnerait} sans doute --- la pince serre les bords extérieurs, non le creux ---, mais \emph{à l'aveugle} : la stratégie ignore que l'\emph{anse} offrirait une prise plus naturelle, et ne sait pas \emph{juger} qu'une prise vaut mieux qu'une autre. Elle saisit sans \emph{comprendre} l'objet. Et --- il faut être juste --- combler ce manque relève autant de la \emph{perception} que de la planification : reconnaître une anse ou un bord comme une prise suppose justement de le \emph{comprendre}, pas seulement de le localiser. C'est là, plus que dans le calcul de l'angle lui-même, que se situe l'axe d'amélioration --- une saisie qui déciderait à partir de la \emph{forme}, de la \emph{pose} et des \emph{prises offertes} par l'objet, ce vers quoi tend précisément une politique apprise (chapitre~\ref{ch:evaluation}).

\paragraph{Le sous-actionnement : une donnée du problème, non un défaut.}
Cet angle piloté par la position n'est pas un caprice de conception : il découle du \emph{sous-actionnement} du bras (chapitre~\ref{ch:planification}). Le SO-101 dispose de cinq articulations motorisées, une de moins que les six degrés de liberté d'une pose \emph{quelconque} : il lui manque le \emph{lacet} de la tête. À une position donnée, on ne peut donc pas régler librement le \emph{cap} de la pince --- seulement son \emph{inclinaison} dans le plan vertical que l'épaule oriente vers l'objet. C'est cette contrainte, et non un choix arbitraire, qui réduit le degré de liberté d'orientation à un simple \emph{angle d'attaque}. Elle explique aussi pourquoi la prise \emph{latérale} --- saisir l'objet par le flanc --- n'a pas été retenue : non qu'elle soit strictement impossible, mais elle est \emph{mal conditionnée} sur cinq articulations (son succès dépend de la posture de départ du solveur), trop peu fiable pour être exploitée (chapitre~\ref{ch:planification}). Un bras à \emph{six} degrés de liberté lèverait la contrainte. Enfin, la branche \emph{de face} ($90^{\circ}$) réservée aux objets \emph{hauts} a bien été vérifiée à la main et fonctionne, mais elle ne figure dans \emph{aucune} campagne : nous n'en tirons donc aucun chiffre, seulement le constat qualitatif qu'elle est atteignable (chapitre~\ref{ch:evaluation}). Cela dessine du même coup une limite de \emph{hauteur} : les prises fiables supposent un objet assez \emph{bas} pour l'approche verticale ou diagonale ; un objet trop haut ne relèverait que de cette branche frontale non évaluée, et sort donc en pratique de l'ensemble des prises éprouvées. Comprendre cette limite mécanique, plutôt que de la combattre, est ce qui a permis de décider \emph{quelles} saisies valaient la peine d'être tentées.

\paragraph{L'open-vocabulaire plutôt qu'un détecteur ré-entraîné.}
Côté perception, le choix discutable est celui du \emph{type} de détecteur appris. Un détecteur à vocabulaire \emph{fermé} comme YOLO~\cite{redmon_yolo_2016}, \emph{ré-entraîné} sur nos objets, serait sans doute plus rapide et plus précis sur ces classes --- mais au prix d'un jeu d'images \emph{annotées} par objet, soit plusieurs jours de collecte et d'annotation, à refaire pour chaque nouvel objet (chapitre~\ref{ch:etat}). Le détecteur à vocabulaire \emph{ouvert} retenu (OWL-ViTv2) supprime ce coût : changer d'objet cible ne demande que de réécrire une \emph{description textuelle}, sans ré-entraînement. Le prix en est une part de précision et une inférence plus lente --- un compromis que nous jugeons raisonnable pour un travail dont l'effort porte sur l'\emph{orchestration} du pipeline, non sur l'optimisation d'un détecteur. Nous n'avons pas essayé d'autres détecteurs à vocabulaire ouvert, tel Grounding-DINO~\cite{liu_grounding_dino_2024}, OWL-ViTv2 s'inscrivant dans le même \emph{écosystème Hugging Face} (outillage et \emph{Hub} de modèles) que les politiques apprises de la seconde approche --- une cohérence qui a pesé dans le choix.

\paragraph{La politique apprise : résolution contre cadence.}
Un dernier arbitrage, propre à la seconde approche. Une politique apprise agit \emph{image par image} : à chaque pas, elle observe puis décide, et sa \emph{réactivité} dépend directement de sa \emph{cadence} d'inférence. Une première version, entraînée en pleine résolution, inférait trop lentement --- de l'ordre de $13$ images par seconde --- pour une boucle réactive, où une politique qui décide à chaque image a besoin d'un flux soutenu pour suivre le mouvement. Ré-entraîner puis exporter le modèle en résolution \emph{réduite} a fait remonter la cadence à environ $24$ images par seconde (le point de contrôle « basse résolution », chapitre~\ref{ch:evaluation}), assez pour agir de façon fluide. Le compromis est explicite : on échange du \emph{détail visuel} contre de la \emph{cadence}. Ce choix, purement pragmatique, a rendu la politique utilisable en boucle fermée --- au prix, peut-être, d'une part de finesse dans la perception, que nous ne cherchons pas à surinterpréter.

\section{Limites assumées : évitement d'obstacles et occlusion}
\label{sec:disc-assumees}

Les limites précédentes, nous les avons \emph{subies} puis corrigées autant que possible. D'autres, au contraire, nous les avons \emph{choisies} : faute de temps et pour faire d'abord fonctionner la saisie, certains objectifs du cahier des charges ont été délibérément mis de côté.

\paragraph{L'évitement d'obstacles : un objectif délibérément écarté.}
Le cahier des charges prévoyait des trajectoires \emph{évitant les obstacles} (objectif~3). Nous l'avons délibérément écarté --- non par oubli, mais par arbitrage : faire d'abord fonctionner la chaîne perception--planification--contrôle sur un objet \emph{seul} nous a paru prioritaire, et un évitement d'obstacles général aurait demandé un effort hors de portée du temps imparti. En conséquence, la trajectoire (l'enchaînement de segments quintiques du chapitre~\ref{ch:controle}) mène à l'objet \emph{sans} vérifier une représentation de la scène : il n'y a pas de \emph{détection de collision}. La scène de test étant mono-objet, sans distracteur, la question ne se posait pas en pratique --- mais c'est bien une limite, et la métrique « collisions évitées » de l'objectif~6 reste, par construction, \emph{vide} (chapitre~\ref{ch:evaluation}). Une planification capable de contourner des obstacles est renvoyée aux perspectives (section~\ref{sec:disc-pistes}).

\paragraph{L'occlusion : un traitement partiel, et pourquoi.}
La sélection de points de vue face aux occlusions (objectif~2) n'est, elle, que \emph{partiellement} traitée. La reconstruction 3D repose sur la \emph{triangulation} stéréo, qui exige que les \emph{deux} caméras fixes voient l'objet ; si l'une d'elles en est empêchée --- un distracteur, un bord de champ ---, le pipeline ne peut plus trianguler et s'\emph{arrête} proprement. Un repli monoculaire par PnP existe bien dans le code (\texttt{enable\_mono\_pnp\_fallback}), mais il reste en pratique \emph{inopérant} sur nos objets, pour deux raisons qui se cumulent (chapitre~\ref{ch:perception}) : nos détecteurs ne livrent qu'un \emph{contour} grossier --- une boîte, non des points d'appui précis ---, et surtout la pose d'un objet quasi \emph{plan} vu de face est géométriquement \emph{ambiguë} pour une seule caméra, quel que soit le solveur. C'est un résultat \emph{négatif} assumé. La piste --- non réalisée --- serait de faire \emph{coopérer} les caméras autrement : qu'une caméra fixe voyant encore l'objet \emph{amorce} la caméra de poignet, ou que cette dernière tente seule une localisation lorsque la paire fixe est aveugle. Nous la présentons comme une perspective, non comme un acquis.

\paragraph{Un périmètre d'objets réduit.}
Enfin, un mot sur le périmètre d'objets. Le plan initial en prévoyait \emph{cinq} ; la \emph{campagne} d'évaluation, elle, s'est concentrée sur le \emph{cube} et, secondairement, sur le \emph{cylindre} (chapitre~\ref{ch:evaluation}). Les autres --- balle, boîte, prisme, et même un \emph{Rubik's cube} visé par le détecteur appris (OWL-ViTv2) --- ont bien été \emph{saisis}, mais de manière \emph{informelle}, sans entrer dans la campagne mesurée, faute de temps. Nous disposons donc d'indices \emph{qualitatifs} que le pipeline s'étend au-delà de ses deux objets mesurés, sans statistiques sur ce périmètre élargi. La conséquence porte sur ce que nous pouvons \emph{affirmer} : la \emph{généralité} du pipeline --- saisir d'autres formes, ou un objet désigné par le langage --- est \emph{étayée} (plusieurs formes pratiquées, le cylindre chiffré) sans être \emph{exhaustivement mesurée}. Le pipeline reste \emph{plus} général que la politique apprise (chapitre~\ref{ch:evaluation}), mais l'étendue exacte de cette généralité, au-delà du cube et du cylindre, demande une campagne élargie pour être quantifiée.

\section{Pistes d'amélioration}
\label{sec:disc-pistes}

Cette mise en perspective ouvre naturellement sur les directions qui prolongeraient le travail. Plusieurs découlent directement des limites ci-dessus.

\paragraph{Perception : segmenter et comprendre l'objet.}
Un module de \emph{segmentation} (tel SAM~2~\cite{ravi_sam2_2024}) fournirait au planificateur un \emph{masque} précis, et non plus une simple boîte englobante : de quoi affiner l'estimation d'orientation, aujourd'hui dépendante d'un contour parfois grossier (section~\ref{sec:disc-limites}), et amorcer une perception des \emph{prises} de l'objet --- repérer l'anse ou le bord d'une tasse au lieu d'en viser le centre (section~\ref{sec:disc-choix}). Un mode \emph{hybride} --- HSV rapide pour le suivi, détecteur appris pour la validation périodique --- concilierait par ailleurs la \emph{cadence} du premier et la \emph{robustesse sémantique} du second.

\paragraph{Planification : contourner les obstacles, décider d'après l'objet.}
Deux directions se rejoignent ici. D'une part, remplacer l'enchaînement de segments quintiques par un \emph{planificateur par échantillonnage} de type RRT~\cite{lavalle_planning_2006}, capable de \emph{contourner} des obstacles --- ce qui adresserait l'objectif~3 délibérément écarté (section~\ref{sec:disc-assumees}). D'autre part, faire décider la \emph{saisie} par la forme et la pose de l'objet, et non par sa seule position : choisir l'angle d'attaque \emph{et} le point de prise d'après l'objet lui-même, en s'appuyant sur la perception enrichie ci-dessus (section~\ref{sec:disc-choix}).

\paragraph{Perception active et occlusion.}
Sélectionner \emph{dynamiquement} le point de vue le moins occulté --- l'objectif~2 du cahier, aujourd'hui seulement effleuré par la coopération stéréo / caméra de poignet --- rendrait la reconstruction robuste aux masquages. Dans le même esprit, une véritable \emph{récupération d'occlusion} ferait coopérer les caméras au lieu d'arrêter le pipeline dès qu'une vue fixe est aveugle (section~\ref{sec:disc-assumees}).

\paragraph{Asservissement visuel continu.}
La correction \texttt{cam\_2} actuelle est \emph{ponctuelle} --- on regarde une fois, on corrige, on descend (\emph{look-and-move}, chapitre~\ref{ch:controle}). La remplacer par un asservissement visuel \emph{continu}, référencé image~\cite{chaumette_visual_2006}, corrigerait la trajectoire à chaque image plutôt qu'une seule fois, au prix d'un détecteur assez rapide pour suivre la cadence.

\paragraph{Matériel : un sixième degré de liberté.}
La limite la plus structurelle n'est ni logicielle ni algorithmique : c'est le \emph{sous-actionnement} du bras (section~\ref{sec:disc-choix}). Un bras à \emph{six} degrés de liberté --- ou une modification du SO-101 ajoutant un \emph{lacet} au poignet --- libérerait l'orientation d'approche et débloquerait les prises \emph{latérales} aujourd'hui mal conditionnées. C'est l'amélioration qui lèverait la contrainte la plus fondamentale du travail --- celle que ni la calibration ni le code ne peuvent contourner.

\paragraph{Un retour d'effort.}
La tenue de l'objet est aujourd'hui vérifiée par un \emph{proxy} --- la charge du moteur de pince (chapitre~\ref{ch:controle}), faute de capteur dédié. Un \emph{capteur tactile} ou d'effort à la pince y substituerait une vraie mesure de contact --- de quoi détecter plus tôt un objet qui glisse ou une fermeture à vide, et \emph{doser} le serrage plutôt que de le tenir constant.

\paragraph{Évaluation et apprentissage.}
Enfin, deux prolongements côté mesure. Élargir la \emph{campagne} formelle à l'ensemble des objets déjà pratiqués (section~\ref{sec:disc-assumees}) chiffrerait la généralité aujourd'hui seulement étayée. Et, du côté de la politique apprise, l'enrichir de démonstrations \emph{plus nombreuses et plus variées} --- d'autres objets, d'autres positions --- élargirait sa distribution, là où elle demeure aujourd'hui \emph{spécialisée} (chapitre~\ref{ch:evaluation}).

\bigskip
\noindent Ces pistes supposent toutes un même socle : une pipeline complète, instrumentée et reproductible. La conclusion qui suit en dresse le bilan d'ensemble et le confronte aux objectifs du cahier des charges.
